\documentclass[11pt]{article}
\usepackage[T1]{fontenc}
\usepackage{lmodern,amsmath,amssymb,amsthm,mathtools}
\usepackage[margin=0.78in]{geometry}
\usepackage{booktabs,longtable,array,microtype,xcolor,enumitem,fancyhdr}
\usepackage{xurl,hyperref}
\hypersetup{colorlinks=true,linkcolor=blue!45!black,urlcolor=blue!55!black,citecolor=blue!45!black,
 pdftitle={Exact Integer Projections for Diophantine Consistency},
 pdfauthor={Research prepared for Vladimir Reshetnikov}}
\setlength{\parskip}{4pt}
\setlength{\parindent}{0pt}
\setlength{\headheight}{14pt}
\pagestyle{fancy}\fancyhf{}
\fancyhead[L]{\small Exact integer projections for Diophantine consistency}
\fancyhead[R]{\small 25 October 2026}\fancyfoot[C]{\thepage}
\allowdisplaybreaks[3]
\newtheorem{theorem}{Theorem}[section]
\newtheorem{lemma}[theorem]{Lemma}
\newtheorem{proposition}[theorem]{Proposition}
\newtheorem{corollary}[theorem]{Corollary}
\newcommand{\Z}{\mathbb Z}
\newcommand{\N}{\mathbb N}
\newcommand{\Con}{\operatorname{Con}}
\newcommand{\ZFC}{\mathrm{ZFC}}
\newcommand{\Ht}{\operatorname{ht}}
\newcommand{\normone}[1]{\left\lVert#1\right\rVert_1}
\title{Exact Integer Projections\\for Diophantine Consistency}
\author{Research report prepared for Vladimir Reshetnikov\\\small Explicit equations, proofs, and independently checked finite artifacts}
\date{25 October 2026}
\begin{document}
\maketitle
\begin{abstract}
We give a completely specified ordinary polynomial in 92 unrestricted
integer variables, of exact total degree 692, whose absence of integer
zeros is provably equivalent in ZFC to syntactic consistency of ZFC.
The new construction is obtained by exact integral projections of an
authenticated direct encoding of the C222 numerical proof calculus.
It uses neither a K215 simulation nor an instance of Jones's universal
polynomial. Every eliminated coordinate has a unique polynomial
restoration with integer coefficients.

The inspected direct predecessor has 94 variables and degree 65,540.
The new equation improves both of those measures, while using 29
residual squares instead of 28. Its expanded coefficient height is
proved strictly below $2^{1655}$; no full expansion of the principal
polynomial is needed or claimed. The complete finite polynomial
definition is printed in this report and supplied as a binary
arithmetic graph. Lower-degree companions include a sextic in 191
integers and a sextic in 193 integers with coefficient height 16,866.

We also close a previously stated transport question for a ten-positive-
variable Pell construction, and give an explicit signed, multiple-base
certificate theorem. The report distinguishes new endpoints, inherited
methods, a late-arriving exact duplicate, and research changes that do
not improve the measured frontier. The new finite algebra is independently
checked; the inherited unbounded arithmetic and logical calibration is
identified rather than represented as a newly compiled proof-assistant
development. Global minimality is not asserted.
\end{abstract}
\tableofcontents
\newpage

\section{The result and the meaning of simplicity}

\subsection{The principal equation}

All unqualified variables in the principal construction range over
$\Z$. Let $S_0,\ldots,S_{127}\in\Z[v_0,\ldots,v_{190}]$ be the 128
literal residuals printed in Appendix~\ref{app:source}. Let $F$ be the
ordered set of 92 free coordinate indices, let $J$ be the set of 29
retained residual indices, and let the 99 acyclic unit definitions be
those printed in Appendix~\ref{app:definition}. These finite data define
polynomials $V_0(\boldsymbol x),\ldots,V_{190}(\boldsymbol x)$ in exactly
the free inputs $\boldsymbol x=(x_0,\ldots,x_{91})$. Define
\begin{equation}\label{eq:principal}
 \boxed{\displaystyle
 P_{92}(\boldsymbol x)=\sum_{r\in J}
       S_r\bigl(V_0(\boldsymbol x),\ldots,V_{190}(\boldsymbol x)\bigr)^2=0.}
\end{equation}
The symbols $V_j$ are polynomial abbreviations. They are not additional
quantified unknowns. The two appendices give a complete coefficient
specification of one fixed ordinary polynomial; a finite acyclic arithmetic
graph provides the same specification in executable form.

\begin{theorem}[Principal endpoint]\label{thm:main}
The polynomial in \eqref{eq:principal} has integer coefficients and exact
total degree $692$. Its expanded coefficient height satisfies
$\Ht(P_{92})<2^{1655}$. Composing the explicit integer projections below
with the retained C222/ZF/ZFC calibration gives
\[
 \ZFC\ \vdash\
 \left(\neg\exists\boldsymbol x\in\Z^{92}\;
                 P_{92}(\boldsymbol x)=0\right)
       \ \longleftrightarrow\ \Con(\ZFC).
\]
Here $\Con(\ZFC)$ means that the specified effective ZFC presentation
has no finite syntactic refutation.
\end{theorem}

The theorem is an equivalence. It does not settle which side holds.
The coefficients, coordinate lists, and arithmetic graph are chosen by
finite construction and do not depend on a consistency assumption.
The formalization boundary is explained in Section~\ref{sec:logic}.

\subsection{The comparison that is actually improved}

Package 152's direct-calculus family contains a 94-integer polynomial
of exact degree 65,540, presented by 28 residual squares, with the
reported expanded-height bound $2^{164870}$~\cite{p152}.
The comparison is therefore:
\begin{center}
\begin{tabular}{lrr}
\toprule
Measure & Inspected 94-variable form & New principal form\\
\midrule
Independent integer variables & 94 & 92\\
Exact total degree & 65,540 & 692\\
Residual squares & 28 & 29\\
Certified expanded-height bound & $<2^{164870}$ & $<2^{1655}$\\
Largest literal scalar in definition & 256 & 46,656\\
\bottomrule
\end{tabular}
\end{center}
The degree is about $94.71$ times smaller. The height entries are bounds,
not two measured exact heights; a comparison of upper bounds is not a
proof that one actual expanded height is smaller. The larger literal
scalar and the extra residual are explicit costs of this particular
choice of source and projection.

This is an improvement within the directly compiled proof-calculus
family. The repository also contains ten-integer constructions through
Sun's reduction, a 37-integer quartic, and nine-positive-variable forms
with different coefficient, degree, or ancestry costs~\cite{catalogue}.
Consequently 92 is not a global minimum for the task. No argument here
establishes a globally optimal polynomial under any total ordering of
variables, degree, coefficient height, or description length.

\subsection{Other completed endpoints}

The lower-degree alternatives are fixed ordinary polynomials in
unrestricted integers. Every residual coefficient is supplied, and every
collected coefficient is supplied for the first six entries below.
\begin{center}\small
\begin{tabular}{rrrrr}
\toprule
Variables&Degree&Residuals&Collected monomials&Exact height\\\midrule
191&6&128&13,074&3,268,176,223\\
193&6&130&25,617&16,866\\
197&6&134&13,552&640,307\\
145&8&82&23,048&604,716,407,814\\
134&12&71&40,248&604,717,858,816\\
119&16&56&90,760&253,229,015,978,920,082,669,568\\
115&24&52&uncollected&unmeasured\\
\bottomrule\end{tabular}\end{center}

For the 115-variable degree-24 form, the full residual lists are supplied
but their sum of squares was not collected. Its exact degree follows
from the maximum residual degree and the sum-of-squares argument below.

The 193-variable sextic is a useful alternative when small literal
coefficients matter: all of its cubic residual coefficients have
absolute value at most $64$, and the fully collected polynomial has
height $16,866$. The 191-variable sextic uses fewer variables and fewer
collected monomials, with larger coefficients. These two points are
different tradeoffs. The 92-variable graph should be used when the
lower dimension and degree-692 profile are the desired objective.

\section{Exact elimination over the integers}\label{sec:elimination}

\subsection{A unit coefficient supplies a polynomial inverse}

\begin{lemma}[Integral monic projection]\label{lem:projection}
Let $\epsilon\in\{-1,1\}$, let $f\in\Z[\boldsymbol x]$, and let
$R_i\in\Z[\boldsymbol x,y]$. Put
\[
 A(\boldsymbol x,y)=(\epsilon y+f(\boldsymbol x))^2+
                    \sum_iR_i(\boldsymbol x,y)^2,
\qquad
 B(\boldsymbol x)=\sum_iR_i(\boldsymbol x,-\epsilon f(\boldsymbol x))^2.
\]
Then the integer zero sets of $A$ and $B$ are in bijection. The map from
the smaller zero set is the integral polynomial map
$\boldsymbol x\mapsto(\boldsymbol x,-\epsilon f(\boldsymbol x))$,
and its inverse is coordinate projection. In fact
\[
 A(\boldsymbol x,-\epsilon f(\boldsymbol x))=B(\boldsymbol x)
\]
as a polynomial identity.
\end{lemma}
\begin{proof}
At an integer zero of $A$, every square is nonnegative, so each residual
is zero. The distinguished row uniquely gives $y=-\epsilon f$.
Conversely this value makes that row zero, and substitution gives the
displayed identity. All restored coordinates are integers because
$\epsilon$ is an integer unit. The two maps are inverse on the stated
zero sets.
\end{proof}

There is no division by an arbitrary coefficient in this lemma. A row
$2y-f=0$, for example, cannot be handled by silently inserting $f/2$.
Nor does the lemma authorize choosing a square root, restricting the
sign of an unrestricted coordinate, or assuming that a modulus is nonzero.
The checker rejects such changes by requiring the literal pivot unit.

A small example explains why degree and dimension move in opposite
directions. Eliminating $y$ from
\[
 (y-ab)^2+(z-yc)^2=0
\]
gives the sextic $(z-abc)^2=0$. One quantified coordinate disappears,
and its multiplication becomes part of the remaining polynomial.
The same operation on an actual compiled proof certificate can remove
many auxiliary coordinates while preserving the complete set of integer
witnesses through explicit maps.

\subsection{Composition and the importance of acyclicity}

Suppose several selected source rows have the form
$\epsilon_v v+f_v=0$, where $f_v$ does not explicitly contain $v$.
It is not enough to inspect these rows separately: the definitions
could contain a cycle. For example, using both $u=v+1$ and $v=u-1$
as though they eliminate both coordinates would leave an unspecified
free parameter.

For the principal construction the complete graph of dependencies is
acyclic. Give the free coordinates their specified values, evaluate
the 99 definitions in topological order, and obtain a unique integer
tuple in the 191-coordinate source. The retained residuals then have
their exact source values. Conversely every source zero satisfies all
of these definitions. Thus the composite map is a polynomial bijection
of integer zero sets, with coordinate projection as inverse.

This graph proof is independent of the heuristic that found the list of
pivots. The release contains two separate checks: sequential replay of
the unit eliminations and a check of the final simultaneous definitions
using topological sorting. Both are exact finite algebra checks.

\subsection{Degree-bounded searches and factored definitions}

The lower-degree profiles were found by considering unit pivots whose
actual substituted residuals remain below a specified degree bound.
The search compares coordinate count, support, and coefficient costs.
It uses exact integer polynomial arithmetic; randomized choices only
decide which eligible pivot to try first. The final certificate fixes
every choice, so verification and reconstruction are deterministic.

For larger degree bounds, fully expanding intermediate polynomials
is unnecessary and can be expensive. The 92-variable definition retains
the finite arithmetic graph. Degree upper bounds propagate through
addition and multiplication, and a separate nonvanishing certificate
proves the exact final degree. There are no variable exponents or
unbounded recursion in this graph. Every node is an integer constant,
a free input, an addition, a subtraction, or a multiplication.

These are applications and improvements of the repository's existing
monic-elimination method~\cite{p152,p159}; the general method is not
claimed as a new discovery. The new work consists of the selected
source combinations, optimized literal endpoints, complete restoration
certificates, and their independent checks.

An initially smaller source need not lead to the best result after a
degree bound is imposed. In the inspected searches, the 220-variable
source exposed a better sextic projection than its more contracted
218-variable sibling. Earlier eliminations can obscure inexpensive
unit pivots or increase the degree of later substitutions. The search
results establish constructions, not lower bounds on what other
coordinate systems could achieve.

\section{The exact degree and coefficient bound}\label{sec:degree}

\subsection{Why the principal degree is exactly 692}

Propagation through the verified arithmetic graph gives maximum
residual degree at most $346$. Therefore $\deg P_{92}\le692$.
The release supplies an explicit integer ray
$x_j=a_jt$, with all 92 values $a_j$ recorded. On that ray, source
residual 61 has coefficient
\[
 [t^{346}]\,S_{61}(V(\boldsymbol a t))
          \equiv 613639\pmod{1000003}.
\]
The residue is nonzero, so the integer coefficient is nonzero.
That residual has exact degree $346$.

One must not now assume that the highest squares cannot cancel modulo
the chosen modulus. The conclusion is instead made over the integers:
if $L_r$ is the leading degree-346 homogeneous part of a residual, then
the degree-692 homogeneous part of $P_{92}$ is $\sum_r L_r^2$.
At least one $L_r$ is nonzero, so this sum is a nonzero real polynomial.
For an entirely coefficient-based proof, take the largest monomial in
a multiplicative monomial order among the $L_r$. Its squared monomial
has a positive sum of squared integer leading coefficients.
Thus $\deg P_{92}=692$.

The same argument proves exact degree $2d$ for each expanded
sum-of-squares profile having a nonzero residual of maximum degree $d$.
Nonzero modular derivative witnesses separately authenticate the use
of every free variable in the principal polynomial.

\subsection{Three different coefficient measures}

For a polynomial $A=\sum_\alpha a_\alpha\boldsymbol x^\alpha$, put
$\normone{A}=\sum_\alpha|a_\alpha|$. Then
\[
 \normone{A+B}\le\normone A+\normone B,
 \qquad \normone{AB}\le\normone A\normone B,
 \qquad \Ht(A)\le\normone A.
\]
Apply these inequalities to every acyclic definition and then to the
sum of residual squares. The resulting explicit positive integer
upper bound has 1655 bits. Hence
\[
                       \Ht(P_{92})<2^{1655}.
\]
This is a rigorous upper bound on the expanded height. The number
$46,656$ is instead the largest absolute literal scalar in the
delivered finite definition. Neither number is a measured exact
expanded height of $P_{92}$, and the full expanded monomial count of
$P_{92}$ is not asserted.

For the expanded companion equations the distinction is simpler:
their complete collected coefficient lists are supplied, so their
reported term counts and heights are exact. This is why the report
labels the principal height as a bound and the companion heights as
measured values.

\section{Logical calibration and proof scope}\label{sec:logic}

The direct route starts with the retained polynomial in 225 unrestricted
integers in package 158~\cite{p158}. Its source arithmetic encodes the
numerical C222 proof calculus, including signed Pell certificates,
finite proof-label conditions, and the associated bounded arithmetic
interfaces. It does not introduce a universal-machine program number
as an unspecified coefficient.

Package 167 supplies exact integer projections to the 220-variable
source used here~\cite{p167}. The new sparse certificate eliminates
29 coordinates to obtain the printed 191-variable sextic. The new
factored certificate then eliminates 99 more coordinates to give
$P_{92}$. The two new stages are unconditional integral polynomial
maps. The separately included 193-variable sextic comes from the
late 217-variable direct quartic, with its own retained parent archive.

The named logical interfaces in the inherited development include
\begin{center}\small
\texttt{KEndpoint.c222\_iff\_zf\_refutation}\\
\texttt{ConstructibleChoice.zf\_refutation\_iff\_zfc\_refutation}.
\end{center}
The supplied calibration audit checks the relation to the effective
ZFC presentation and the internalizability of the argument in ZFC.
That argument uses semantic model constructions, completeness,
well-founded-part constructions, and the constructible inner model
for the Choice bridge. This report does not replace that argument by
a claim that the entire calibration is merely finite arithmetic or
has been proved in PA. No external transitivity or standardness
assumption is added to the statement of Theorem~\ref{thm:main}.

For the new projection layer, only the elementary unit-elimination
argument is needed. If the inherited source polynomial is called $Q$,
the explicitly checked restoration gives
\[
 (\exists\boldsymbol x\in\Z^{92}\ P_{92}=0)
 \ \longleftrightarrow\
 (\exists\boldsymbol y\ Q=0)
 \ \longleftrightarrow\
 \text{there is a finite syntactic ZFC refutation}.
\]
Negating the equivalence yields the stated consistency assertion.

The new package supplies ordinary mathematical proofs formalizable
in ZFC and exact executable checks of the finite polynomials and maps.
It does not supply a newly compiled end-to-end Lean or Rocq proof,
nor a serialized object-language derivation inside ZFC. The original
archives and named source interfaces are retained so that these
dependencies can be audited. The finite checkers do not search for a
zero, prove consistency, or replace an unbounded arithmetic theorem
by testing small examples.

\section{A squared transport coordinate in the compact family}

\label{sec:compact-companion}

The direct-calculus constructions and the following compact companion
optimize different quantities.  This section gives two literal polynomials
in ten \emph{strictly positive} integer variables, each of exact degree
$232$ and $85$ binary arithmetic operations.  The trace parent has
$87$ operations at degree $232$; the inverse-Lucas parent has $85$
operations at degree $304$.  The new squared-coordinate variant also
reduces the degree of one carrier numerator from $57$ to $21$.
The total degree remains $232$.

\subsubsection*{Late comparison and attribution}
At the later read-only snapshot
\texttt{0d75e561a8cefc8c30fc4224605a8984c810ad1c}, the incoming
Pooled Powers Separated Radicals archive already contains exactly the
same $85$-gate $bV$ polynomial: all rows, all ten inputs and all four fixed
coefficient values agree.  The $85/232$ metric point is therefore not
claimed as new here.  That equation is a reproduced baseline.  The
distinct contribution is the single replacement $L=bV\mapsto b^2$
and its soundness/completeness proof, which lowers the carrier degree
without changing total degree or operation count.  The exact comparison
is recorded in the companion's \texttt{LATE\_COMPARISON.json}.

The two parents are packages \texttt{163-trace-radical} and
\texttt{164-unsquared-factors-linear-selectors} of the pinned ProveIt
snapshot.  The former explicitly left open the squared-coordinate
transport argument.  Its difficulty was a correction growing with a
negative rank.  Evaluating the inverse correction at a positive Pell
coordinate removes that difficulty.  This proves the particular new
variant below; it does not assert worldwide novelty or global minimality.
This compact component retains the native/K215 ancestry of its fixed
coefficient tuple.

\subsection{Complete polynomial}

The quantified inputs are
\[
 J,F,\alpha,i,T,w,s,n,\eta,\zeta\in\mathbb Z_{>0}.
\]
The four fixed integers $\beta,K_0,M_C,M_F$ are the complete indexed
BFS943/radix67 tuple, reconstructed from the supplied finite recipe.
In particular $\beta=2^{590069}-1$; the largest port has $592884$ bits.
All four actual ports, rather than arbitrary integers with the same
elementary bounds, are needed for the logical calibration.
Define polynomial abbreviations
\begin{align*}
 h&=\beta J,&q&=h+1,& C&=q-F-\alpha,\\
 X&=qw,&Y&=qs,&k&=\eta+\zeta,\\
 c&=kY+\eta,&A&=Y(X+1)+1,&\Delta&=A^2-1,\\
 H&=4A-5,&v&=XY^2,&\ell&=vk,\\
 E&=1+\ell(\ell+k),&&&
\end{align*}
and
\begin{align*}
 R&=[h(q-F)+\alpha](q^2-1)+(M_C+qM_F)J,\\
 B&=(K_0+w)C-F,\\
 U&=ic^2,\qquad a=\Delta U^2E,\qquad F_0=1+a,\qquad G=EF_0,\\
 f&=1+2a,\qquad Q=4\Delta aF_0,\\
 b&=cT-R,\qquad V=fb-c,\qquad K=X(3c-2X)+2,\\
 e_6&=6(k-R-1)-\ell(k^2-4),\qquad D_6=6q^6,\\
 M&=hD_6H,\\
 Z&=hD_6[H(n-LB)+K]-H[D_6B+h(e_6-6q^4s)] .
\end{align*}
The principal weight is $L=b^2$; a companion uses $L=bV$.
The equation is
\begin{equation}
 \boxed{\;Q\bigl(Z^2-M^2GV^2\bigr)+\bigl(M^2G-Z^2\bigr)=0.\;}
 \label{eq:compact232}
\end{equation}
Every symbol on the right is a polynomial abbreviation.  There are no
additional quantified variables, variable exponents, recurrence operators,
or uncomputed coefficients.  The literal circuit computes
$\ell=(XY)(kY)$ and $f=a+F_0$, thereby sharing existing operations.

Composed with the named inherited native/refutation calibration, the
ordinary mathematical theorem is
\[
 \mathrm{ZFC}\vdash
 \left(\neg\exists\boldsymbol{x}\in\mathbb Z_{>0}^{10}\;
 P(\boldsymbol{x})=0\right)
 \longleftrightarrow\operatorname{Con}(\mathrm{ZFC}).
\]
The domain is essential.  If nonnegative variables are desired, substitute
$x_j=u_j+1$ for each of the ten inputs.  This leaves the dimension and
total degree unchanged, though the displayed operation count refers to
the positive-input circuit.

\subsection{The new transport proof}

For the proof only, put
\[
 e_*={e_6\over6},\quad \rho={e_*\over q^6},\quad
 \sigma={s\over q^2},\quad t={B\over h},\quad
 \kappa={K\over H},\quad W=1+hL,\quad z={Z\over M}.
\]
No integrality of these quotients is initially assumed.  The exact
rational identity is
\begin{equation}
 z=n+\kappa+\sigma-Wt-\rho.\label{eq:compact-carrier}
\end{equation}

The trace parent's root extraction and low-index exclusions depend on
$Q,V,G$ and the final norm, not on $L$ or $e_6$.
At a positive zero they therefore still give
\[
 G=g^2,\quad z=\varepsilon gy,\quad
 QV^2-(Q-1)y^2=1,\qquad g,y>0,\quad\varepsilon\in\{-1,1\},
\]
with $z,y$ integral, and
\[
 b\ne0,\quad |b|>8f^3-\tfrac12,\quad
 g<f/2,\quad y<3|V|/2,\quad c<f/2.
\]
These statements are established before any sign assumption on $R$ or $B$.
In particular, since $b$ is a nonzero integer,
\[
 |V|\le f|b|+c<\tfrac32 f|b|,
 \qquad
 gy<\tfrac98 f^2|b|<\tfrac12 b^2.
\]
The last strict inequality follows from
$|b|>8f^3-1/2>(9/4)f^2$ for $f\ge1$.
Consequently $L=b^2$ satisfies
\[
 L>0,\qquad gy<L/2,\qquad |b|/q^6<L/2.
\]

The inverse correction has the decisive sign:
\[
 {vk(k^2-4)\over6}\ge0,
 \qquad e_*\le k-R-1<b,
\]
because $k\ge2$ and $b-(k-R-1)=cT-k+1>0$.
This holds for every integer $R$, including negative ranks.  Thus
$\rho<|b|/q^6$.  If $B\le-1$, then $Wt<-L$; using the positive
$n,\kappa,\sigma$ in \eqref{eq:compact-carrier} gives
\[
 n=\varepsilon gy-\kappa-\sigma+Wt+\rho
   <L/2-L+|b|/q^6<0,
\]
a contradiction.  Therefore $B\ge0$ and then $C>0$.
This is the new unbounded sign argument.  It is not inferred from
bounded testing or from a formal substitution into the forward-correction
proof.

\subsection{The remaining arithmetic order}

Once $C>0$, the inherited mask bound gives $3q<R<q^4-2$.
The unsquared trace obeys $f\equiv1\pmod c$, so the two auxiliary
residues remain
$V\equiv-c\pmod f$ and $V\equiv-R\pmod c$.
The retained signed-index theorem authenticates the main Pell rank as $R$.

For clarity, the inverse denominator argument must retain its order.
The first Pell recurrence has
\[
 k=2m+8vC_m+32v^2Z_m,\qquad
 C_m=\binom{m+1}{3},\quad Z_m\in\mathbb N.
\]
Putting $d=2m-R-1$, $A_m=C_m+4vZ_m$ and
$S_m=k^2+2mk+4m^2-4$ gives the exact identity
\[
 3(e_*-d)=4v^2(24Z_m-A_mS_m).
\]
The carrier gives $q^6\mid(e_*-q^4s)$, hence $q^4\mid e_*$.
Since $v=q^3ws^2$, prime valuations in the displayed identity give
$q^4\mid(e_*-d)$ even at the prime $3$.
The bounds $|d|<R<q^4$ then force $d=0$.

It would be premature to cancel the remaining factor $3$ and infer
$q^6\mid e_*$.  Instead, the carrier first gives $H\mid6hK$.
The retained uniform ratio estimate and whole-product projection force
$X=2^R$.  Since $q\mid X$, the radix $q$ is a power of two.  Only now
does the identity imply $q^6\mid e_*$.  The carrier then gives
$q^2\mid s$ and $h\mid B$, supplying the exact native outer endpoint.
The parent's indexed native/refutation calibration completes soundness.

Conversely, a finite refutation gives the canonical outer witnesses and
the trace parent's positive auxiliary Pell pair.  Its actual positive
integer input is $T=(V+c+Rf)/(cf)$.
The inverse error can be negative, but the inherited explicit bounds give
\[
 |\rho|<vk^3/3<g/3,\qquad \kappa+\sigma<g/2.
\]
For either positive weight choose
\[
 n=gy+(1+hL)t+\rho-\kappa-\sigma.
\]
This is an integer and is strictly positive because $y\ge1$, $t\ge0$
and $g>\kappa+\sigma+|\rho|$.  The carrier now has $z=gy$, proving
completeness for all ten positive inputs.

\subsection{Exact complexity and scope}

Let $q_0=\beta J$, $k_0=\eta+\zeta$ and
$B_0=k_0sT-q_0^2(q_0-F)$.  The complete highest homogeneous component is
\[
 [P]_{232}=-9\,2^{10}q_0^{100}w^{26}s^{58}i^{10}k_0^{32}B_0^2.
\]
It is nonzero at the actual $\beta>0$.
For $L=b^2$, the degrees of $(Q,M,G,V,Z)$ are
$(80,11,50,40,21)$.  The term $QZ^2$ has degree $122$, while
$-QM^2GV^2$ has degree $232$ and remains the unique dominant summand.
For $L=bV$, $\deg Z=57$ and $\deg(QZ^2)=194$.

The reproducible finite checker covers every live gate by $36$ exact
free-polynomial cuts per graph, checks the complete $202$-monomial highest
component and all four coefficient payloads, and rejects $22$ deliberately
corrupted operators per graph.  Normal and optimized Python runs agree.
The finite recipe reconstructs every fixed integer and matches its pinned
raw SHA-256 identity.  These are algebraic artifact checks, separate from
the unbounded arguments and inherited logical calibration.
No new end-to-end proof-assistant development is claimed.

\subsection{A smaller, separately calibrated coefficient tuple}

Both $85$-operation, degree-$232$ forms also admit the complete inherited
\emph{complemented-label radix47} specialization.  All four fixed ports
change together, while every free-port gate remains identical.  This
separate specialization has
\[
 \beta=2^{413929}-1,\qquad
 \text{largest port}=415904\text{ bits},\qquad
 \text{total raw bytes}=207214.
\]
The radix67 tuple has $295388$ raw bytes, so the inherited smaller tuple
saves $88174$ bytes in this matched storage measure.  Its coefficient
reduction is credited to its native predecessor; the present contribution
is its explicit composition with the new $85/232$ polynomials.

This is not a coefficient-storage record.  The late incoming comparison
already supplies a separately calibrated one-anchor7309 tuple using
$172011$ raw bytes.  Its bV graph also duplicates the $85/232$ baseline,
as disclosed above.  The present distinct carrier change remains $L=b^2$.

This is the whole complemented-label tuple with seven shared marker weights,
not the conditional47 layout and not the radix67 formulas with a substituted
radix.  The complete ordinary carry/controller proof, its native bridge,
the $3297$-coefficient array, all $51$ native target identities and the
literal native P66 polynomial are retained and authenticated against the
bridge's original source hashes.  The finite recipe reconstructs every
port and checks the actual controller at prime $1012691$, its exact orders
and its valuation-one lift.

The new trace and inverse-correction arguments are uniform in the four
ports subject to the matched scalar bounds, mask phases and population
identities.  They recover the same outer endpoint with $W_{\rm base}=w$.
The whole47 bridge then reconstructs fresh source witnesses at
$A_{\rm source}=Y(X+1)+2$ and applies its own complemented-label decoder.
Conversely its own controller supplies the canonical positive outer data;
the new bound $g>\kappa+\sigma+|\rho|$ proves the constructed $n$ positive.
Thus both directions use this actual whole tuple.  No old67 decoder or
pointwise identification of the two specialized polynomials is assumed.

The complete alternate coefficient recipe is supplied separately with the
radix47 component.  Its independent verifier checks the whole binding,
all live gates, full degree-$232$ leading forms, and actual corruptions.
It carries the same explicitly stated ordinary-proof and formalization
boundary as the main compact companion.


\section{Oriented rational error certificates}\label{sec:pooling}

The incoming reciprocal-pooling construction combines two consecutive
bases~\cite{p168}. Its final research paragraph already suggests
oriented weights for several signed bases. The following result closes
the size and zero-base details for five blocks and is implemented in
the supporting 209-variable degree-24 equation. That endpoint is not
competitive with all the exact-projection profiles above; its value
here is the completed arithmetic certificate and its explicit compiler.

\begin{theorem}[Five signed blocks]\label{thm:pooling}
Suppose $W\ge258$, $|A|>W^W$, and there are at most five integer bases
$n_j$, each satisfying $n_j^2<W$, and at most 17 requests altogether.
For each request let $m_i,k_i\in\Z$ satisfy $m_i^2,k_i^2<W$, and let $(X_i,Y_i)$ be
the canonical signed Pell coordinates at parameter $A$ and index $k_i$.
Set
\[
 T_j=2An_j-n_j^2-1,\quad D_i=m_i-X_i+AY_i,
\]
\[
 H_j=\sum_{i\in I_j}Y_i(m_i-X_i),\qquad
 E_j=\sum_{i\in I_j}(D_i^2-Y_i^2)-2n_jH_j.
\]
Then the existence of an integer $z$ satisfying
\begin{equation}\label{eq:oriented}
 \sum_j n_jE_j\prod_{\ell\ne j}T_\ell
                  -z\prod_jT_j=0
\end{equation}
implies $m_i=n_j^{k_i}$ for every block with $n_j\ne0$.
At nonunit bases it also implies $k_i\ge0$. Conversely, these true
powers in all nonzero-base blocks imply the existence of such a $z$.
Zero-base blocks impose no condition in this theorem.
\end{theorem}
\begin{proof}
For a nonzero base, $|T_j|>W^W$ and
$\operatorname{sign}(T_j)=\operatorname{sign}(An_j)$.
For a zero base, $T_j=-1$, so none of the denominators is zero.
Put $\delta_i=D_i-n_jY_i$. Direct expansion gives
\[
 \sum_{i\in I_j}\delta_i^2=E_j-T_j\sum_{i\in I_j}Y_i^2.
\]
Consequently \eqref{eq:oriented} makes
$\sum_j n_j\sum_i\delta_i^2/T_j$ an integer.

Let $b=\lfloor\sqrt{W-1}\rfloor\ge16$ and take the product only over
nonzero bases in
$C=\prod_{n_j\ne0}n_j^b$. For each request in such a block define
\[
 d_i=Cm_i-(C/n_j^b)n_j^{b+k_i}\in\Z.
\]
The exponents are nonnegative and the indicated quotient is an integer.
The canonical Pell congruence gives $C\delta_i\equiv d_i\pmod{T_j}$.
It remains valid for negative indices because $\gcd(n_j,T_j)=1$.
It follows that $\sum_j n_j\sum_i d_i^2/T_j$ is an integer.

There are at most five nonzero factors in $C$. Thus
$|d_i|\le2b^{6b}$ and
\[
 \sum_j |n_j|\sum_{i\in I_j}d_i^2
 \le68b^{12b+1}<W^{6b+2}\le W^W.
\]
Here $b^2<W$, $68b<W^2$, and $6b+2\le W$ for $b\ge16$.
The integer rational sum has absolute value below one. Moreover every
nonzero weight $n_j/T_j$ has the same sign, namely the sign of $A$.
Hence that integer is zero and every $d_i$ is zero. Dividing in $\mathbb Q$
by the nonzero integer $C$ proves the power identities. Negative indices
at a nonunit base cannot give an integer output.

Conversely true powers give $T_j\mid\delta_i$, whence $T_j\mid E_j$.
The value $z=\sum_{n_j\ne0}n_jE_j/T_j$ is an integer and satisfies
\eqref{eq:oriented}. The terms from zero bases are zero.
\end{proof}

The actual five bases are $F^2+1$, $B$, $BH$, $q$, and $q+1$.
Their nonvanishing is recovered in a checked order: the retained norm
forces $M\ne0$; the relation $Bv=M$ forces $B\ne0$; oddness gives
$H\ne0$. The first three blocks can now be interpreted. The retained
rounding bound gives $|B|\ge2$, and the exact powers give
$q=(BH)^x\ne0,-1$. Therefore the two remaining blocks are also covered.
This avoids assuming the later radix positivity in order to authenticate
the earlier power conditions.

The supporting compiler removes ten private certificate coordinates
and introduces one integer $z$ in the protected 218-variable source.
The resulting 209-variable polynomial has 145 quadratic residuals and
one residual of degree 12, hence exact total degree 24. The independent
checker authenticates the protected row basis and every replacement.
No expanded support or coefficient-height record is claimed for it.

\section{A sound semantic alternative that does not improve the frontier}

The supplementary factored-radix construction replaces the request
$q=(BH)^x$ by
\[
                  N=B^x,\qquad P=H^x,\qquad q=NP.
\]
Its changed frame uses 35 semantic atom squares. The outputs and indices
retain squared bounds by $W$, while the two bases $q,q+1$ need only the
weaker bounds $|q|,|q+1|<W$. The same-base certificate remains sound:
with $b=\lfloor\sqrt{W-1}\rfloor$, the cleared error bound becomes
\[
 60W^{4b}<W^{4b+1}\le W^W,
 \qquad |2An-n^2-1|\ge2|A|-W^2>W^W.
\]
The full proof checks the signed cases and the ordering of witness choices.
In particular the new frame contains only semantic coordinates chosen
before the large Pell parameter, not private Pell witnesses whose size
depends on that parameter.

The literal source has 224 variables, 161 residuals, and degree six;
five rechecked radical projections give a 219-variable source. The
subsequent searches reached 193-variable sextics, so this changed
semantics has not improved the frozen 191-variable sextic frontier.
It is included as completed research with that limitation, not counted
as an unachieved 190-variable result.

\section{Verification, provenance, and integration}\label{sec:verification}

\subsection{What was checked}

The principal component binds its 191-variable source to the exact
220-to-191 endpoint, checks all 99 unit definitions and their acyclic
dependencies, checks all retained residuals, and supplies a literal
binary arithmetic graph. It independently verifies the degree witness,
the height bound, and actual occurrence of the free coordinates.
The integer restoration can be evaluated at arbitrary supplied tuples.
The resulting value of the source sum of squares agrees exactly with
the value of the projected polynomial.

For the sparse companions, a separately implemented verifier uses
exponent maps rather than the producer's repeated-index representation.
It substitutes through every remaining row, checks the full endpoint,
renumbers coordinates, and compares every collected coefficient where
an expansion is supplied. Deliberate pivot, restoration, and endpoint
corruptions must be rejected, including under optimized Python.
The compact and rational-pooling components have their own exact
identity and source-binding checks. Repeated small numerical tests
are not offered as proof of an unbounded assertion.

The arithmetic-parent replays and the exact scope of each fresh run
are recorded in the verification receipts. Included historical receipts
are not represented as fresh runs of their entire transitive proof
development. Independent reviews in this report are separate tasks
within an AI-assisted research workflow, not outside human refereeing.

\subsection{Immutable comparison boundaries}

The initial inspection was pinned to commit
\texttt{35fa5f6bc659e83f2b2034126c4bcf8e023fbe1c}.
The repository changed during the work. A late inspection at
\texttt{0d75e561a8cefc8c30fc4224605a8984c810ad1c}
covered newly placed packages 166--168 and three new incoming archives.
It found direct quartics with 216 or 217 variables and an exact duplicate
of the compact $bV$ 85-operation degree-232 graph. Those findings are
included in the comparisons and source ledger.

In particular the compact 85/232 metric is not claimed as a new record,
and the smaller historical radix47 coefficient tuple is not described
as the current storage record. The distinct compact result here is
the proved $b^2$ carrier. The new direct projection endpoints are
compared with the identified sources, rather than an unspecified
moving repository head.

\subsection{Using the archive}

The principal equation, its binary graph, source, certificates,
evaluator, and independent checker are under \texttt{principal92/}.
The lower-degree ordinary equations are under \texttt{equations/}.
All these definitions use literal integers and finite operations.
The original archives retained under \texttt{inherited/} preserve
the arithmetic and logical proof dependencies and their notices.

From the extracted release directory, the principal commands are:
\begin{verbatim}
python3 -B principal92/verify.py
python3 -B verify.py
python3 -B -O verify.py
python3 -B verify.py --components
\end{verbatim}
Each successful command checks a finite artifact. It is not a procedure
for deciding whether the polynomial has an integer zero. The README
gives exact evaluation and restoration commands, the report build
instructions, and optional parent-replay commands.

For additive integration into ProveIt, keep the literal equation and
the source-binding certificates together. A useful formalization target
is the generic integral monic-projection theorem plus a checked
interpreter for the finite graph. The inherited C222 and ZF/ZFC
calibration remains a separately identifiable dependency.

\appendix
\section{Complete finite definition of the 92-variable polynomial}
\label{app:definition}

The next data, together with the source rows in Appendix~\ref{app:source},
specify \eqref{eq:principal} without a coefficient oracle or an
unspecified program parameter. Let the increasing free list be
$F=(F_0,\ldots,F_{91})$. Set $V_{F_j}=x_j$. For each listed triple
$(r,v,\epsilon)$ define the polynomial
\[
 f_{r,v}(\boldsymbol u)=-\epsilon
       \bigl(S_r(\boldsymbol u)-\epsilon u_v\bigr),\qquad
 V_v=f_{r,v}(V_0,\ldots,V_{190}).
\]
The displayed source row contains $u_v$ only as the term
$\epsilon u_v$. The table is ordered so that dependencies on other
defined coordinates precede their use. It therefore specifies
polynomial abbreviations, not simultaneous existential equations.
Finally square and sum the rows indexed by $J$.
\subsection*{Free coordinates and retained residuals}
\begin{align*}
F={}&(4,5,6,7,8,9,10,11,12,13,14,15,17,18,20,42,\\
&{}46,50,78,80,81,82,84,85,87,88,90,91,93,94,96,97,\\
&{}98,100,101,103,105,107,108,110,111,113,114,117,121,123,124,125,\\
&{}126,127,129,130,132,133,135,136,138,139,141,142,144,145,147,148,\\
&{}150,151,154,156,157,159,160,162,163,166,167,171,172,173,174,175,\\
&{}178,179,180,181,182,183,184,185,186,187,189,190)\\
J={}&(13,27,31,35,61,77,79,84,86,88,90,92,94,96,98,100,\\
&{}102,104,106,108,112,114,119,121,122,123,124,126,127)\\
\end{align*}
\subsection*{The 99 unit definitions in dependency order}
\begin{longtable}{rrr@{\qquad}rrr@{\qquad}rrr@{\qquad}rrr}
\toprule
$r$&$v$&$\epsilon$&$r$&$v$&$\epsilon$&$r$&$v$&$\epsilon$&$r$&$v$&$\epsilon$\\\midrule\endhead
110&0&1&4&19&1&5&21&1&6&22&1\\
7&23&1&18&33&1&32&47&1&36&51&1\\
47&62&1&120&104&1&115&168&1&2&1&1\\
1&2&-1&9&25&1&11&27&1&12&28&1\\
14&29&1&16&31&1&22&37&1&33&48&1\\
38&53&1&48&63&1&65&86&-1&66&89&-1\\
67&92&-1&68&95&-1&69&99&-1&70&102&-1\\
71&106&-1&72&109&-1&73&112&-1&74&115&-1\\
116&116&1&0&3&1&8&24&1&15&30&1\\
17&32&1&19&34&1&21&36&1&23&38&1\\
34&49&1&37&52&1&40&55&1&51&65&1\\
78&122&1&83&128&1&85&131&1&87&134&1\\
89&137&1&91&140&1&93&143&1&95&146&1\\
97&149&1&99&152&1&81&153&-1&101&155&1\\
103&158&1&105&161&1&107&164&1&111&176&1\\
125&188&-1&3&16&1&10&26&1&20&35&1\\
24&39&1&28&43&1&39&54&1&41&56&-1\\
42&57&-1&43&58&-1&44&59&-1&45&60&-1\\
46&61&-1&50&66&1&52&67&1&80&120&-1\\
117&169&1&25&40&1&29&44&1&49&64&1\\
53&68&1&60&75&1&62&76&1&63&77&1\\
109&165&1&26&41&1&30&45&1&54&69&1\\
82&79&1&118&170&1&55&70&1&64&83&-1\\
75&118&1&113&177&1&56&71&1&76&119&1\\
57&72&1&58&73&1&59&74&1&&&\\
\bottomrule\end{longtable}


\section{All 128 literal source residuals}\label{app:source}

These are the complete residuals $S_r(\boldsymbol v)$ of the
191-variable sextic source, indexed from zero. Every coefficient
and exponent below is literal. The source equation is
$\sum_{r=0}^{127}S_r(\boldsymbol v)^2=0$.
% Complete literal source rows; generated from source191.json.
\begingroup\small\setlength{\jot}{3pt}
\begin{align*}
S_{0}={}&-v_{2}^{2}+v_{3}\\
S_{1}={}&v_{0}v_{5}v_{7}-v_{0}v_{7}-v_{2}-v_{5}v_{6}+v_{5}v_{7}\\
S_{2}={}&-1-v_{0}v_{5}v_{7}+v_{0}v_{7}+v_{1}-v_{5}v_{7}\\
S_{3}={}&1+v_{3}-v_{4}v_{18}+v_{16}\\
S_{4}={}&-1+2v_{8}-v_{8}^{2}-2v_{8}v_{9}+2v_{9}-v_{9}^{2}+v_{19}\\
S_{5}={}&-2+2v_{8}-v_{8}^{2}-4v_{8}v_{20}+4v_{20}-4v_{20}^{2}+v_{21}\\
S_{6}={}&-2-v_{9}^{2}-2v_{9}v_{10}-v_{10}^{2}+v_{22}\\
S_{7}={}&-1+2v_{8}-v_{8}^{2}-2v_{8}v_{10}+2v_{10}-v_{10}^{2}+v_{23}\\
S_{8}={}&19-26v_{8}-29v_{8}^{2}-34v_{8}v_{9}-36v_{8}v_{10}-24v_{8}v_{19}-24v_{8}v_{23}\\
&{}+41v_{9}+38v_{9}^{2}-48v_{9}v_{10}+16v_{9}v_{22}-28v_{9}v_{23}+36v_{10}+70v_{10}^{2}\\
&{}-32v_{10}v_{19}+32v_{10}v_{22}-13v_{19}-8v_{19}v_{23}-v_{20}+40v_{22}-14v_{23}\\
&{}+v_{24}-v_{27}\\
S_{9}={}&-6+16v_{10}-16v_{10}^{2}-16v_{10}v_{22}-16v_{10}v_{23}+8v_{22}-4v_{22}^{2}-8v_{22}v_{23}\\
&{}+8v_{23}-4v_{23}^{2}+v_{25}\\
S_{10}={}&6+36v_{8}-48v_{8}v_{10}-48v_{8}v_{23}-24v_{8}v_{25}+48v_{9}-64v_{9}v_{10}\\
&{}-64v_{9}v_{23}-32v_{9}v_{25}+7v_{10}+9v_{10}^{2}-16v_{10}v_{19}+8v_{10}v_{23}\\
&{}+4v_{10}v_{25}+12v_{19}-16v_{19}v_{23}-8v_{19}v_{25}-v_{20}+8v_{23}+4v_{25}\\
&{}+v_{26}-v_{36}\\
S_{11}={}&96-29v_{8}-37v_{8}^{2}-24v_{8}v_{9}-112v_{8}v_{10}-4v_{8}v_{22}-24v_{8}v_{23}\\
&{}-39v_{9}+17v_{9}^{2}-4v_{9}v_{10}-24v_{9}v_{22}-4v_{9}v_{23}-80v_{10}-74v_{10}^{2}\\
&{}-32v_{10}v_{22}-32v_{10}v_{23}-66v_{22}-4v_{22}^{2}-8v_{23}-4v_{23}^{2}+v_{27}\\
S_{12}={}&-12+16v_{8}-4v_{8}^{2}-8v_{8}v_{10}-8v_{8}v_{19}-8v_{8}v_{22}+16v_{10}-4v_{10}^{2}\\
&{}-8v_{10}v_{19}-8v_{10}v_{22}+16v_{19}-4v_{19}^{2}-8v_{19}v_{22}+16v_{22}-4v_{22}^{2}\\
&{}+v_{28}\\
S_{13}={}&-4-30v_{8}-8v_{8}^{2}-64v_{8}v_{9}+16v_{8}v_{19}-48v_{9}-64v_{9}^{2}-6v_{19}\\
&{}+8v_{19}^{2}+2v_{20}-v_{31}\\
S_{14}={}&-82-414v_{8}-529v_{8}^{2}-1472v_{8}v_{9}-368v_{8}v_{19}+92v_{8}v_{20}-576v_{9}\\
&{}-1024v_{9}^{2}-512v_{9}v_{19}+128v_{9}v_{20}-144v_{19}-64v_{19}^{2}+32v_{19}v_{20}\\
&{}+36v_{20}-4v_{20}^{2}+v_{29}\\
S_{15}={}&-1453-10332v_{8}-17139v_{8}^{2}-46656v_{8}v_{9}-11664v_{8}v_{19}\\
&{}+2916v_{8}v_{20}+744v_{8}v_{29}-14016v_{9}-31744v_{9}^{2}-15872v_{9}v_{19}\\
&{}+3968v_{9}v_{20}+1024v_{9}v_{29}-3504v_{19}-1984v_{19}^{2}+992v_{19}v_{20}\\
&{}+256v_{19}v_{29}+876v_{20}-124v_{20}^{2}-64v_{20}v_{29}+225v_{29}-4v_{29}^{2}+v_{30}\\
S_{16}={}&-1-28v_{8}^{2}-32v_{8}v_{9}-40v_{8}v_{19}+8v_{9}-32v_{9}v_{19}-4v_{19}-12v_{19}^{2}\\
&{}-v_{20}+v_{31}\\
S_{17}={}&49-34v_{9}^{2}-44v_{9}v_{10}-24v_{9}v_{22}-32v_{9}v_{31}+6v_{10}-10v_{10}^{2}\\
&{}-24v_{10}v_{22}-24v_{10}v_{31}-32v_{22}-8v_{22}v_{31}+4v_{31}-4v_{31}^{2}+v_{32}\\
&{}-v_{37}\\
S_{18}={}&-v_{10}^{2}-2v_{10}v_{11}-v_{11}^{2}+v_{33}\\
S_{19}={}&-26+14v_{8}+v_{8}^{2}-8v_{8}v_{9}+8v_{8}v_{10}-8v_{8}v_{22}-8v_{8}v_{23}\\
&{}-4v_{8}v_{33}+8v_{9}^{2}+16v_{9}v_{22}-8v_{9}v_{23}-8v_{10}+2v_{10}^{2}-8v_{10}v_{33}\\
&{}+8v_{22}^{2}-8v_{22}v_{23}+18v_{23}-4v_{23}^{2}+4v_{33}-4v_{33}^{2}+v_{34}-3v_{37}\\
S_{20}={}&4v_{8}-18v_{8}^{2}-40v_{8}v_{10}-32v_{8}v_{23}-40v_{8}v_{33}-16v_{8}v_{34}\\
&{}+4v_{10}-24v_{10}^{2}-32v_{10}v_{23}-48v_{10}v_{33}-16v_{10}v_{34}-16v_{23}^{2}\\
&{}-32v_{23}v_{33}-16v_{23}v_{34}+4v_{33}-24v_{33}^{2}-16v_{33}v_{34}-4v_{34}^{2}+v_{35}\\
S_{21}={}&v_{10}-25v_{10}^{2}-40v_{10}v_{23}-20v_{10}v_{25}-16v_{23}^{2}-16v_{23}v_{25}-4v_{25}^{2}\\
&{}+v_{36}\\
S_{22}={}&-1+8v_{9}-4v_{9}^{2}-8v_{9}v_{22}+8v_{22}-4v_{22}^{2}+v_{37}\\
S_{23}={}&-1-16v_{9}^{2}-8v_{9}v_{10}-32v_{9}v_{22}-16v_{9}v_{37}-v_{10}^{2}-8v_{10}v_{22}\\
&{}-4v_{10}v_{37}-16v_{22}^{2}-16v_{22}v_{37}-4v_{37}^{2}+v_{38}\\
S_{24}={}&3+v_{9}-2v_{9}v_{10}-4v_{9}v_{38}-3v_{10}^{2}-8v_{10}v_{38}-v_{22}-4v_{38}^{2}\\
&{}+v_{39}\\
S_{25}={}&-1-4v_{39}-4v_{39}^{2}+v_{40}\\
S_{26}={}&-1-v_{10}^{2}-8v_{10}v_{39}-4v_{10}v_{40}-16v_{39}^{2}-16v_{39}v_{40}-4v_{40}^{2}+v_{41}\\
S_{27}={}&4+v_{9}-2v_{9}v_{10}-4v_{9}v_{41}-3v_{10}^{2}-8v_{10}v_{41}-v_{22}-4v_{41}^{2}\\
&{}+v_{42}\\
S_{28}={}&-4v_{36}^{2}-8v_{36}v_{42}-4v_{42}^{2}+v_{43}\\
S_{29}={}&-9+7v_{8}-v_{8}^{2}-8v_{8}v_{36}-4v_{8}v_{43}+24v_{36}-16v_{36}^{2}-16v_{36}v_{43}\\
&{}+12v_{43}-4v_{43}^{2}+v_{44}\\
S_{30}={}&-1-4v_{44}-4v_{44}^{2}+v_{45}\\
S_{31}={}&-v_{10}^{2}-8v_{10}v_{44}-4v_{10}v_{45}-16v_{44}^{2}-16v_{44}v_{45}-4v_{45}^{2}+v_{46}\\
S_{32}={}&-v_{8}-v_{8}^{2}-4v_{8}v_{10}-4v_{8}v_{46}-4v_{10}-4v_{10}^{2}-8v_{10}v_{46}-4v_{46}\\
&{}-4v_{46}^{2}+v_{47}\\
S_{33}={}&-10-12v_{47}-4v_{47}^{2}+v_{48}\\
S_{34}={}&2-4v_{9}-16v_{9}v_{47}-8v_{9}v_{48}-4v_{22}-16v_{22}v_{47}-8v_{22}v_{48}\\
&{}-v_{37}+8v_{47}-16v_{47}^{2}-16v_{47}v_{48}+4v_{48}-4v_{48}^{2}+v_{49}\\
S_{35}={}&-10+v_{9}-25v_{9}^{2}-40v_{9}v_{22}-20v_{9}v_{49}-16v_{22}^{2}-16v_{22}v_{49}\\
&{}-12v_{42}-4v_{42}^{2}-4v_{49}^{2}+v_{50}\\
S_{36}={}&21+24v_{42}-8v_{42}^{2}-16v_{42}v_{50}-2v_{50}-4v_{50}^{2}+v_{51}\\
S_{37}={}&1+4v_{8}-4v_{8}^{2}-8v_{8}v_{19}-8v_{8}v_{53}+v_{10}-v_{10}^{2}-4v_{10}v_{51}\\
&{}+4v_{19}-4v_{19}^{2}-8v_{19}v_{53}-4v_{51}^{2}+v_{52}+4v_{53}-4v_{53}^{2}\\
S_{38}={}&-4v_{8}-4v_{8}^{2}-4v_{8}v_{10}-8v_{8}v_{19}-v_{10}-v_{10}^{2}-4v_{10}v_{19}-4v_{19}\\
&{}-4v_{19}^{2}+v_{53}\\
S_{39}={}&4v_{8}-24v_{8}^{2}-48v_{8}v_{19}-16v_{8}v_{52}-16v_{8}v_{53}+4v_{19}-24v_{19}^{2}\\
&{}-16v_{19}v_{52}-16v_{19}v_{53}-4v_{52}^{2}+8v_{53}-8v_{53}^{2}+v_{54}\\
S_{40}={}&-v_{2}v_{13}+v_{55}\\
S_{41}={}&-6v_{8}^{2}v_{55}-120v_{8}v_{20}v_{55}-48v_{8}v_{21}v_{55}-120v_{8}v_{55}+v_{12}\\
&{}-v_{13}-248v_{20}^{2}v_{55}-128v_{20}v_{21}v_{55}-232v_{20}v_{55}-8v_{21}^{2}v_{55}\\
&{}-62v_{21}v_{55}+134v_{55}-v_{56}\\
S_{42}={}&-4v_{8}^{2}v_{55}-16v_{8}v_{9}v_{55}-16v_{8}v_{22}v_{55}+28v_{8}v_{55}+8v_{9}^{2}v_{55}\\
&{}+16v_{9}v_{22}v_{55}+v_{12}-v_{13}+8v_{22}^{2}v_{55}-8v_{24}^{2}v_{55}-6v_{37}v_{55}\\
&{}-30v_{55}-v_{57}\\
S_{43}={}&-4v_{8}^{2}v_{55}-48v_{8}v_{9}v_{55}-64v_{8}v_{10}v_{55}-16v_{8}v_{22}v_{55}\\
&{}-20v_{8}v_{55}-36v_{9}^{2}v_{55}-168v_{9}v_{10}v_{55}-96v_{9}v_{22}v_{55}\\
&{}-144v_{9}v_{55}-148v_{10}^{2}v_{55}-128v_{10}v_{22}v_{55}-192v_{10}v_{55}+v_{12}\\
&{}-v_{13}-16v_{22}^{2}v_{55}-156v_{22}v_{55}-8v_{27}^{2}v_{55}+180v_{55}-v_{58}\\
S_{44}={}&-96v_{8}^{2}v_{55}-256v_{8}v_{10}v_{55}-128v_{8}v_{22}v_{55}-64v_{8}v_{23}v_{55}\\
&{}-64v_{8}v_{28}v_{55}-112v_{8}v_{55}-104v_{10}^{2}v_{55}-96v_{10}v_{22}v_{55}\\
&{}-16v_{10}v_{23}v_{55}-80v_{10}v_{28}v_{55}-180v_{10}v_{55}+v_{12}-v_{13}\\
&{}-16v_{22}^{2}v_{55}-32v_{22}v_{28}v_{55}-48v_{22}v_{55}+8v_{23}^{2}v_{55}-16v_{23}v_{28}v_{55}\\
&{}-76v_{23}v_{55}-4v_{25}v_{55}-8v_{28}^{2}v_{55}-8v_{28}v_{55}+56v_{55}-v_{59}\\
S_{45}={}&-2v_{9}^{2}v_{55}-8v_{9}v_{30}v_{55}+v_{12}-v_{13}-8v_{30}^{2}v_{55}-4v_{30}v_{55}\\
&{}-v_{60}\\
S_{46}={}&-128v_{8}^{2}v_{55}-192v_{8}v_{10}v_{55}-64v_{8}v_{23}v_{55}-64v_{8}v_{32}v_{55}\\
&{}+16v_{8}v_{55}-72v_{10}^{2}v_{55}-48v_{10}v_{23}v_{55}-48v_{10}v_{32}v_{55}\\
&{}+12v_{10}v_{55}+v_{12}-v_{13}-8v_{23}^{2}v_{55}-16v_{23}v_{32}v_{55}+4v_{23}v_{55}\\
&{}-8v_{32}^{2}v_{55}+8v_{32}v_{55}-v_{61}\\
S_{47}={}&-v_{8}v_{9}+v_{8}v_{10}+v_{9}+v_{9}v_{10}-v_{10}-v_{10}^{2}+v_{62}\\
S_{48}={}&1+2v_{11}+v_{11}^{2}-v_{62}^{2}+v_{63}\\
S_{49}={}&-8v_{54}^{2}v_{55}-12v_{54}v_{55}-2v_{55}+v_{64}\\
S_{50}={}&2v_{2}v_{8}v_{65}+2v_{2}v_{19}v_{65}-2v_{2}v_{65}-v_{9}v_{55}-v_{13}+v_{65}\\
&{}+v_{66}\\
S_{51}={}&-v_{2}v_{8}v_{12}+v_{2}v_{12}-v_{12}+v_{65}\\
S_{52}={}&v_{2}v_{9}v_{12}-6v_{8}v_{55}-8v_{9}v_{55}+v_{12}-v_{13}-2v_{19}v_{55}\\
&{}+v_{67}\\
S_{53}={}&-v_{12}v_{56}v_{57}+v_{13}v_{56}v_{57}+v_{68}\\
S_{54}={}&-v_{12}v_{68}+v_{13}v_{68}+2v_{26}v_{55}v_{68}+6v_{55}v_{68}+v_{69}\\
S_{55}={}&-v_{58}v_{59}v_{69}+v_{70}\\
S_{56}={}&-v_{60}v_{61}v_{70}+v_{71}\\
S_{57}={}&-v_{12}v_{71}+v_{13}v_{71}+2v_{35}v_{55}v_{71}+4v_{55}v_{71}+v_{72}\\
S_{58}={}&-v_{66}v_{67}v_{72}+v_{73}\\
S_{59}={}&-v_{73}^{2}+v_{74}\\
S_{60}={}&-v_{5}v_{14}-v_{5}v_{14}v_{16}+v_{14}v_{16}+v_{75}\\
S_{61}={}&-2v_{1}^{2}+2v_{1}v_{12}+2v_{1}v_{13}+2v_{1}v_{75}+2v_{1}v_{76}-v_{12}^{2}\\
&{}-v_{12}^{2}v_{74}+2v_{12}v_{13}v_{74}+2v_{12}v_{64}v_{74}-2v_{12}v_{75}-v_{13}^{2}\\
&{}-v_{13}^{2}v_{74}-2v_{13}v_{64}v_{74}-2v_{13}v_{76}+v_{17}v_{78}-v_{63}^{2}v_{74}\\
&{}-v_{64}^{2}v_{74}-v_{75}^{2}-v_{76}^{2}\\
S_{62}={}&-2v_{5}v_{15}-v_{5}v_{15}v_{16}+v_{15}+v_{15}v_{16}+v_{76}\\
S_{63}={}&330-217v_{8}-238v_{8}^{2}-64v_{8}v_{9}-512v_{8}v_{10}-120v_{8}v_{20}\\
&{}-48v_{8}v_{21}-160v_{8}v_{22}-128v_{8}v_{23}-64v_{8}v_{28}-64v_{8}v_{32}\\
&{}-153v_{9}-30v_{9}^{2}-168v_{9}v_{10}-80v_{9}v_{22}-8v_{9}v_{30}-360v_{10}\\
&{}-324v_{10}^{2}-224v_{10}v_{22}-64v_{10}v_{23}-80v_{10}v_{28}-48v_{10}v_{32}-4v_{19}\\
&{}-232v_{20}-248v_{20}^{2}-128v_{20}v_{21}-62v_{21}-8v_{21}^{2}-204v_{22}-24v_{22}^{2}\\
&{}-32v_{22}v_{28}-72v_{23}-16v_{23}v_{28}-16v_{23}v_{32}-8v_{24}^{2}-4v_{25}-2v_{26}\\
&{}-8v_{27}^{2}-8v_{28}-8v_{28}^{2}-4v_{30}-8v_{30}^{2}+8v_{32}-8v_{32}^{2}-2v_{35}-6v_{37}\\
&{}-12v_{54}-8v_{54}^{2}+v_{77}\\
S_{64}={}&-v_{4}v_{104}-v_{79}v_{82}v_{104}+v_{80}v_{179}+v_{82}v_{104}-v_{83}\\
S_{65}={}&-v_{17}v_{85}v_{104}+v_{84}v_{179}-v_{86}\\
S_{66}={}&-v_{17}v_{88}v_{104}+v_{87}v_{179}-v_{89}\\
S_{67}={}&-v_{17}v_{91}v_{104}+v_{90}v_{179}-v_{92}\\
S_{68}={}&-v_{17}v_{94}v_{104}+v_{93}v_{179}-v_{95}\\
S_{69}={}&-v_{17}v_{98}v_{104}+v_{96}v_{180}-v_{99}\\
S_{70}={}&-v_{17}v_{101}v_{104}+v_{100}v_{180}-v_{102}\\
S_{71}={}&-v_{17}v_{104}v_{105}+v_{103}-v_{106}\\
S_{72}={}&-v_{17}v_{104}v_{108}+v_{107}-v_{109}\\
S_{73}={}&-v_{17}v_{104}v_{111}+v_{110}v_{179}-v_{112}\\
S_{74}={}&-v_{17}v_{104}v_{114}+v_{113}v_{179}-v_{115}\\
S_{75}={}&4+2v_{0}+3v_{0}^{2}-6v_{0}^{2}v_{117}+2v_{0}v_{3}-4v_{0}v_{3}v_{117}-4v_{0}v_{117}\\
&{}+v_{1}^{2}-2v_{1}^{2}v_{117}+2v_{3}+2v_{3}^{2}-4v_{3}^{2}v_{117}-4v_{3}v_{117}+v_{4}^{2}\\
&{}-2v_{4}^{2}v_{117}+2v_{5}^{2}-4v_{5}^{2}v_{117}-2v_{8}+v_{8}^{2}-2v_{8}^{2}v_{117}\\
&{}+4v_{8}v_{117}-2v_{9}v_{10}+4v_{9}v_{10}v_{117}+v_{11}^{2}-2v_{11}^{2}v_{117}+v_{12}^{2}\\
&{}-2v_{12}^{2}v_{117}+v_{13}^{2}-2v_{13}^{2}v_{117}+v_{14}^{2}-2v_{14}^{2}v_{117}+v_{15}^{2}\\
&{}-2v_{15}^{2}v_{117}+v_{16}^{2}-2v_{16}^{2}v_{117}+v_{17}^{2}-2v_{17}^{2}v_{117}+v_{22}\\
&{}-2v_{22}v_{117}+v_{77}^{2}-2v_{77}^{2}v_{117}-2v_{79}+2v_{79}^{2}-4v_{79}^{2}v_{117}\\
&{}+4v_{79}v_{117}+v_{80}^{2}-2v_{80}^{2}v_{117}+v_{81}^{2}-2v_{81}^{2}v_{117}+v_{84}^{2}\\
&{}-2v_{84}^{2}v_{117}+v_{87}^{2}-2v_{87}^{2}v_{117}+v_{90}^{2}-2v_{90}^{2}v_{117}+v_{93}^{2}\\
&{}-2v_{93}^{2}v_{117}+v_{96}^{2}-2v_{96}^{2}v_{117}+v_{97}^{2}-2v_{97}^{2}v_{117}+v_{100}^{2}\\
&{}-2v_{100}^{2}v_{117}+v_{103}^{2}-2v_{103}^{2}v_{117}+v_{104}^{2}-2v_{104}^{2}v_{117}+v_{107}^{2}\\
&{}-2v_{107}^{2}v_{117}+v_{110}^{2}-2v_{110}^{2}v_{117}+v_{113}^{2}-2v_{113}^{2}v_{117}-8v_{117}\\
&{}-2v_{117}v_{165}^{2}+v_{118}+v_{165}^{2}\\
S_{76}={}&-3v_{0}^{2}v_{118}-2v_{0}v_{3}v_{118}-2v_{0}v_{118}-v_{1}^{2}v_{118}-2v_{3}^{2}v_{118}\\
&{}-2v_{3}v_{118}-v_{4}^{2}v_{118}-2v_{5}^{2}v_{118}-v_{8}^{2}v_{118}+2v_{8}v_{118}\\
&{}+2v_{9}v_{10}v_{118}-v_{11}^{2}v_{118}-v_{12}^{2}v_{118}-v_{13}^{2}v_{118}-v_{14}^{2}v_{118}\\
&{}-v_{15}^{2}v_{118}-v_{16}^{2}v_{118}-v_{17}^{2}v_{118}-v_{22}v_{118}-v_{77}^{2}v_{118}\\
&{}-2v_{79}^{2}v_{118}+2v_{79}v_{118}-v_{80}^{2}v_{118}-v_{81}^{2}v_{118}-v_{84}^{2}v_{118}\\
&{}-v_{87}^{2}v_{118}-v_{90}^{2}v_{118}-v_{93}^{2}v_{118}-v_{96}^{2}v_{118}-v_{97}^{2}v_{118}\\
&{}-v_{100}^{2}v_{118}-v_{103}^{2}v_{118}-v_{104}^{2}v_{118}-v_{107}^{2}v_{118}-v_{110}^{2}v_{118}\\
&{}-v_{113}^{2}v_{118}-5v_{118}-v_{118}v_{165}^{2}+v_{119}\\
S_{77}={}&2v_{116}+v_{116}^{2}+v_{118}^{2}-v_{119}^{2}\\
S_{78}={}&-v_{116}v_{121}-v_{121}+v_{122}\\
S_{79}={}&-1+v_{120}^{2}+2v_{120}v_{122}+v_{121}^{2}\\
S_{80}={}&-v_{3}^{2}v_{123}+2v_{3}v_{116}v_{123}-v_{3}v_{121}+v_{5}+2v_{116}v_{123}\\
&{}-v_{120}-v_{121}\\
S_{81}={}&-v_{4}^{2}v_{124}+2v_{4}v_{116}v_{124}+2v_{4}v_{124}-v_{4}v_{154}+v_{17}-v_{124}\\
&{}-v_{153}\\
S_{82}={}&-1+v_{79}+2v_{116}v_{125}v_{165}-v_{125}+2v_{125}v_{165}-v_{125}v_{165}^{2}\\
&{}-v_{153}-v_{154}v_{165}\\
S_{83}={}&-v_{116}v_{127}-v_{127}+v_{128}\\
S_{84}={}&-1+v_{126}^{2}+2v_{126}v_{128}+v_{127}^{2}\\
S_{85}={}&1+v_{116}-v_{116}v_{130}-v_{130}+v_{131}\\
S_{86}={}&v_{129}^{2}+2v_{129}v_{131}-2v_{130}+v_{130}^{2}\\
S_{87}={}&-v_{116}v_{133}-v_{133}+v_{134}\\
S_{88}={}&-1+v_{132}^{2}+2v_{132}v_{134}+v_{133}^{2}\\
S_{89}={}&-v_{116}v_{136}-v_{136}+v_{137}\\
S_{90}={}&-1+v_{135}^{2}+2v_{135}v_{137}+v_{136}^{2}\\
S_{91}={}&-v_{116}v_{139}-v_{139}+v_{140}\\
S_{92}={}&-1+v_{138}^{2}+2v_{138}v_{140}+v_{139}^{2}\\
S_{93}={}&-v_{116}v_{142}-v_{142}+v_{143}\\
S_{94}={}&-1+v_{141}^{2}+2v_{141}v_{143}+v_{142}^{2}\\
S_{95}={}&-v_{116}v_{145}-v_{145}+v_{146}\\
S_{96}={}&-1+v_{144}^{2}+2v_{144}v_{146}+v_{145}^{2}\\
S_{97}={}&-v_{116}v_{148}-v_{148}+v_{149}\\
S_{98}={}&-1+v_{147}^{2}+2v_{147}v_{149}+v_{148}^{2}\\
S_{99}={}&-v_{116}v_{151}-v_{151}+v_{152}\\
S_{100}={}&-1+v_{150}^{2}+2v_{150}v_{152}+v_{151}^{2}\\
S_{101}={}&1-v_{116}v_{154}-v_{154}+v_{155}\\
S_{102}={}&-1+2v_{153}+v_{153}^{2}+2v_{153}v_{155}+v_{154}^{2}\\
S_{103}={}&-v_{116}v_{157}-v_{157}+v_{158}\\
S_{104}={}&-1+v_{156}^{2}+2v_{156}v_{158}+v_{157}^{2}\\
S_{105}={}&-v_{116}v_{160}-v_{160}+v_{161}\\
S_{106}={}&-1+v_{159}^{2}+2v_{159}v_{161}+v_{160}^{2}\\
S_{107}={}&-v_{116}v_{163}-v_{163}+v_{164}\\
S_{108}={}&-1+v_{162}^{2}+2v_{162}v_{164}+v_{163}^{2}\\
S_{109}={}&v_{4}-2v_{4}v_{8}-2v_{4}v_{9}-2v_{4}v_{10}-2v_{4}v_{11}-2v_{4}v_{12}\\
&{}-2v_{4}v_{13}-2v_{4}v_{14}-2v_{4}v_{15}-2v_{4}v_{16}+v_{165}\\
S_{110}={}&-16+v_{0}-v_{166}^{2}\\
S_{111}={}&-v_{116}v_{168}-v_{168}+v_{176}\\
S_{112}={}&-1+v_{167}^{2}+v_{168}^{2}-v_{176}^{2}\\
S_{113}={}&-v_{116}v_{170}-v_{167}v_{170}v_{173}-v_{170}+v_{177}\\
S_{114}={}&-1+v_{169}^{2}+v_{170}^{2}-v_{177}^{2}\\
S_{115}={}&-v_{163}^{2}-2v_{163}^{2}v_{172}+v_{168}\\
S_{116}={}&v_{116}+v_{167}v_{173}-4v_{168}-8v_{168}v_{171}\\
S_{117}={}&-v_{162}-v_{164}-v_{167}v_{174}+v_{169}\\
S_{118}={}&-v_{77}-4v_{163}v_{175}+v_{170}\\
S_{119}={}&26+2v_{0}^{2}+2v_{0}v_{3}-2v_{0}v_{121}-2v_{0}v_{127}+v_{1}^{2}-2v_{1}v_{154}\\
&{}+v_{3}^{2}-2v_{3}v_{127}+v_{5}^{2}-2v_{5}v_{145}+v_{8}^{2}-2v_{8}v_{130}+v_{9}^{2}\\
&{}-2v_{9}v_{133}+v_{10}^{2}-2v_{10}v_{136}+v_{11}^{2}-2v_{11}v_{139}+v_{12}^{2}\\
&{}-2v_{12}v_{142}+v_{13}^{2}-2v_{13}v_{148}+v_{14}^{2}-2v_{14}v_{151}+v_{15}^{2}\\
&{}-2v_{15}v_{157}+v_{16}^{2}-2v_{16}v_{160}+v_{77}^{2}-2v_{77}v_{163}-v_{116}v_{178}\\
&{}-2v_{120}-2v_{122}+v_{122}^{2}-2v_{126}-2v_{128}+v_{128}^{2}-2v_{129}+2v_{130}\\
&{}-2v_{131}+v_{131}^{2}-2v_{132}-2v_{134}+v_{134}^{2}-2v_{135}-2v_{137}+v_{137}^{2}\\
&{}-2v_{138}-2v_{140}+v_{140}^{2}-2v_{141}-2v_{143}+v_{143}^{2}-2v_{144}-2v_{146}\\
&{}+v_{146}^{2}-2v_{147}-2v_{149}+v_{149}^{2}-2v_{150}-2v_{152}+v_{152}^{2}-2v_{153}\\
&{}+v_{155}^{2}-2v_{156}-2v_{158}+v_{158}^{2}-2v_{159}-2v_{161}+v_{161}^{2}-2v_{162}\\
&{}-2v_{164}+v_{164}^{2}\\
S_{120}={}&-v_{81}v_{179}+v_{104}\\
S_{121}={}&-v_{97}v_{180}+v_{104}\\
S_{122}={}&-2-v_{77}^{2}-v_{121}^{2}-v_{127}^{2}+2v_{130}-v_{130}^{2}-v_{133}^{2}-v_{136}^{2}-v_{139}^{2}\\
&{}-v_{142}^{2}-v_{145}^{2}-v_{148}^{2}-v_{151}^{2}-v_{154}^{2}-v_{157}^{2}-v_{160}^{2}+v_{163}^{2}\\
&{}-v_{181}^{2}-v_{182}^{2}-v_{183}-v_{183}^{2}\\
S_{123}={}&-1-4v_{83}^{2}-4v_{86}^{2}-4v_{89}^{2}-4v_{92}^{2}-4v_{95}^{2}-4v_{99}^{2}-4v_{102}^{2}+v_{104}^{2}-4v_{106}^{2}\\
&{}-4v_{109}^{2}-4v_{112}^{2}-4v_{115}^{2}-v_{184}^{2}-v_{185}^{2}-v_{186}-v_{186}^{2}-4v_{188}^{2}\\
S_{124}={}&-v_{4}v_{187}+v_{104}\\
S_{125}={}&v_{2}v_{104}-v_{5}v_{187}-v_{188}\\
S_{126}={}&1+v_{77}^{2}+v_{79}^{2}v_{189}-2v_{79}v_{80}v_{127}+2v_{79}v_{84}-2v_{79}v_{84}v_{130}\\
&{}-2v_{79}v_{87}v_{133}-2v_{79}v_{90}v_{136}-2v_{79}v_{93}v_{139}-2v_{79}v_{96}v_{142}\\
&{}-2v_{79}v_{100}v_{148}-2v_{79}v_{103}v_{151}-2v_{79}v_{107}v_{157}\\
&{}-2v_{79}v_{110}v_{160}-2v_{79}v_{113}v_{163}-2v_{79}v_{116}v_{189}\\
&{}+2v_{79}v_{126}v_{127}+2v_{79}v_{127}v_{128}-2v_{79}v_{129}+2v_{79}v_{129}v_{130}\\
&{}+2v_{79}v_{130}v_{131}-2v_{79}v_{131}+2v_{79}v_{132}v_{133}+2v_{79}v_{133}v_{134}\\
&{}+2v_{79}v_{135}v_{136}+2v_{79}v_{136}v_{137}+2v_{79}v_{138}v_{139}\\
&{}+2v_{79}v_{139}v_{140}+2v_{79}v_{141}v_{142}+2v_{79}v_{142}v_{143}\\
&{}+2v_{79}v_{147}v_{148}+2v_{79}v_{148}v_{149}+2v_{79}v_{150}v_{151}\\
&{}+2v_{79}v_{151}v_{152}+2v_{79}v_{156}v_{157}+2v_{79}v_{157}v_{158}\\
&{}+2v_{79}v_{159}v_{160}+2v_{79}v_{160}v_{161}+2v_{79}v_{162}v_{163}\\
&{}+2v_{79}v_{163}v_{164}-2v_{79}v_{189}+v_{80}^{2}-2v_{80}v_{126}+v_{84}^{2}-2v_{84}v_{129}\\
&{}+v_{87}^{2}-2v_{87}v_{132}+v_{90}^{2}-2v_{90}v_{135}+v_{93}^{2}-2v_{93}v_{138}+v_{96}^{2}\\
&{}-2v_{96}v_{141}+v_{100}^{2}-2v_{100}v_{147}+v_{103}^{2}-2v_{103}v_{150}+v_{107}^{2}\\
&{}-2v_{107}v_{156}+v_{110}^{2}-2v_{110}v_{159}+v_{113}^{2}-2v_{113}v_{162}+v_{121}^{2}\\
&{}+v_{126}^{2}+v_{129}^{2}+v_{132}^{2}+v_{135}^{2}+v_{138}^{2}+v_{141}^{2}+v_{145}^{2}+v_{147}^{2}\\
&{}+v_{150}^{2}+v_{154}^{2}+v_{156}^{2}+v_{159}^{2}+v_{162}^{2}-2v_{163}^{2}+v_{181}^{2}+v_{182}^{2}\\
&{}+v_{183}+v_{183}^{2}+v_{189}\\
S_{127}={}&v_{79}^{2}v_{190}-2v_{79}v_{81}v_{121}-2v_{79}v_{97}v_{145}-2v_{79}v_{104}v_{154}\\
&{}-2v_{79}v_{116}v_{190}+2v_{79}v_{120}v_{121}+2v_{79}v_{121}v_{122}\\
&{}+2v_{79}v_{144}v_{145}+2v_{79}v_{145}v_{146}+2v_{79}v_{153}v_{154}+2v_{79}v_{154}\\
&{}+2v_{79}v_{154}v_{155}-4v_{79}v_{190}+v_{81}^{2}-2v_{81}v_{120}+2v_{81}v_{121}+v_{97}^{2}\\
&{}-2v_{97}v_{144}+2v_{97}v_{145}+v_{104}^{2}-2v_{104}v_{153}+2v_{104}v_{154}\\
&{}+2v_{116}v_{190}+v_{120}^{2}-2v_{120}v_{121}-v_{121}^{2}-2v_{121}v_{122}+v_{144}^{2}\\
&{}-2v_{144}v_{145}-v_{145}^{2}-2v_{145}v_{146}+v_{153}^{2}-2v_{153}v_{154}-2v_{154}\\
&{}-v_{154}^{2}-2v_{154}v_{155}+4v_{190}\\
\end{align*}
\endgroup


\begin{thebibliography}{20}
\bibitem{p152} ProveIt, package 152, \emph{Lucas corrections and elimination
of arithmetic witnesses}. Pinned commit 35fa5f6bc659e83f2b2034126c4bcf8e023fbe1c.
\url{https://github.com/VladimirReshetnikov/ProveIt/tree/35fa5f6bc659e83f2b2034126c4bcf8e023fbe1c/Computability/DiophantineConsistency/con-zfc-diophantine-forms/152-lucas-corrections}.
\bibitem{p158} ProveIt, package 158, \emph{Pooling exact power certificates
for Diophantine consistency}. Literal parent P225, retained with original
proofs and notices. \url{https://github.com/VladimirReshetnikov/ProveIt/tree/35fa5f6bc659e83f2b2034126c4bcf8e023fbe1c/Computability/DiophantineConsistency/con-zfc-diophantine-forms/158-pooled-powers}.
\bibitem{p159} ProveIt, package 159, \emph{Quadratic projection 234}.
The general integral radical and unit-projection methods are inherited.
\url{https://github.com/VladimirReshetnikov/ProveIt/tree/35fa5f6bc659e83f2b2034126c4bcf8e023fbe1c/Computability/DiophantineConsistency/con-zfc-diophantine-forms/159-quadratic-projection-234}.
\bibitem{p163} ProveIt, package 163, \emph{Lower degree and fewer witnesses
for Diophantine consistency}; especially compact proof Sections 3--9.
\url{https://github.com/VladimirReshetnikov/ProveIt/tree/35fa5f6bc659e83f2b2034126c4bcf8e023fbe1c/Computability/DiophantineConsistency/con-zfc-diophantine-forms/163-trace-radical}.
\bibitem{p164} ProveIt, package 164, \emph{Unsquared factors and linear
selectors}; inverse Lucas correction and complemented-label radix47 bridge.
\url{https://github.com/VladimirReshetnikov/ProveIt/tree/35fa5f6bc659e83f2b2034126c4bcf8e023fbe1c/Computability/DiophantineConsistency/con-zfc-diophantine-forms/164-unsquared-factors-linear-selectors}.
\bibitem{p167} ProveIt, package 167, \emph{Two smaller direct quartics for
ZFC consistency}. Original Integral Projection archive retained verbatim.
\url{https://github.com/VladimirReshetnikov/ProveIt/tree/0d75e561a8cefc8c30fc4224605a8984c810ad1c/Computability/DiophantineConsistency/con-zfc-diophantine-forms/167-integral-projection}.
\bibitem{p168} ProveIt, package 168, \emph{Reciprocal pooling and integral
projections for Con(ZFC)}. Original archive retained verbatim.
\url{https://github.com/VladimirReshetnikov/ProveIt/tree/0d75e561a8cefc8c30fc4224605a8984c810ad1c/Computability/DiophantineConsistency/con-zfc-diophantine-forms/168-reciprocal-pooling}.
\bibitem{late} ProveIt incoming archives at commit
0d75e561a8cefc8c30fc4224605a8984c810ad1c:
\emph{Pooled Powers Separated Radicals}, \emph{Pooled Projection},
and \emph{Quotient Carriers}, 25 October 2026. Exact archive hashes and
inspection scope are retained in the late-comparison ledger.
\url{https://github.com/VladimirReshetnikov/ProveIt/tree/0d75e561a8cefc8c30fc4224605a8984c810ad1c/docs/incoming}.
\bibitem{catalogue} ProveIt, \emph{Con-ZFC Diophantine forms catalogue},
inspected snapshot retained under \texttt{sources/}; the catalogue is a
source index, not a new independent proof of every listed endpoint.
\url{https://github.com/VladimirReshetnikov/ProveIt/tree/35fa5f6bc659e83f2b2034126c4bcf8e023fbe1c/Computability/DiophantineConsistency/con-zfc-diophantine-forms}.
\end{thebibliography}
\end{document}
